Misalkan
\mathl{P \in N}{} dan
\mathl{P \in U}{} suatu daerah peta dari $M$. Misalkan
\maabbdisp {\alpha} { U } {V
} {}
suatu peta dengan himpunan terbuka
\mathl{V \subseteq H}{,}
\mathl{H=\R_{\geq 0} \times \R^{n-1}}{.} Peta ini menginduksi suatu homeomorfisme
\maabbdisp {} {N \cap U} { \alpha (N \cap U)
} {.}
Peta-peta inilah yang kita gunakan untuk $N$. Sifat difeomorfisme dari pergantian peta untuk $M$ langsung diwariskan oleh pergantian peta untuk $N$. Misalkan
\mathl{P \in N \cap \partial M}{.} Di bawah suatu peta, $P$ kemudian dipetakan ke suatu titik batas dari $H$. Karena peta-peta untuk $N$ merupakan pembatasan dari peta-peta semacam itu, sifat ini tetap berlaku. Sebaliknya, jika
\mathl{P \in \partial N}{}, maka tentu saja, di satu pihak,
\mathl{P \in N \subseteq M}{.} Di pihak lain, di sekitar $P$ terdapat suatu peta untuk $N$ yang diperoleh dengan membatasi suatu peta untuk $M$, dan dalam peta tersebut $P$ dipetakan ke suatu titik batas dari $H$.
